\subsection{交错凯莱代数}

命题 2. 设 A 为一个环，F 为一个凯莱 A-代数，e 为其单位元， \(s: x \mapsto \bar{x}\) 为其共轭，N: F \(\to\) A 为其凯莱范数（§ 3，第 4 号）。

\begin{enumerate}
\def\labelenumi{(\roman{enumi})}
\item
  \(\mathbf{F}\) 为交错代数的必要充分条件是：对 \(x, y\) （ \(\mathbf{F}\) 中元素的每个有序对），有 \(x^2y = x(xy)\) 。
\item
  若 \(\mathbf{F}\) 是交错的，则 \(\mathbf{N}(xy) = \mathbf{N}(x)\mathbf{N}(y)\) 对所有 \(x,y\) （属于 \(\mathbf{F}\) ）成立。
\item
  假设 \(\mathbf{F}\) 是交错的。元素 \(x \in \mathbf{F}\) 可逆的必要充分条件是 \(\mathbf{N}(x)\) 在 \(\mathbf{Ae}\) 中可逆；此时 \(x\) 的逆元是唯一的，且等于 \(\mathbf{N}(x)^{-1}\bar{x}\) ；把它记作 \(x^{-1}\) ，
\end{enumerate}

\[
x ^ {- 1} (x y) = x \left(x ^ {- 1} y\right) = y
\]

对所有 \(y \in \mathbf{F}\) 成立。

条件 \(x^{2}y = x(xy)\) 对于 F 是交错代数显然是必要的（第 1 号，命题 1）。反之，若它对 F 的元素的每个有序对都成立，将其应用于 \(\bar{x}\) 和 \(\bar{y}\) ，就得到 \(\bar{x}^{2}\bar{y} = \bar{x}(\bar{x}\bar{y})\) ；对这一关系式应用共轭 s，我们得到 \(yx^{2} = (yx)x\) ，因此第 1 号命题 1 的条件 (c) 得以满足。

显然 \(a(e, x, y) = 0\) 对所有 \(x, y\) （属于 \(\mathbf{F}\) ）成立。若 \(\mathbf{F}\) 是交错的，则由命题 1（第 1 号）推出，子代数 \(\mathbf{G}\) （ \(\mathbf{F}\) 中由 \(e, x\) 和 \(y\) 生成的子代数）是结合的。由于 \(\bar{x} = -x + \mathrm{T}(x) \in -x + \mathrm{Ae}, \bar{x} \in \mathbf{G}\) ，类似地 \(\bar{y} \in \mathbf{G}\) 。于是 \(\mathrm{N}(xy) = (xy)(\overline{xy}) = xy.\bar{y}.\bar{x} = \mathrm{N}(y)x\bar{x} = \mathrm{N}(y)\mathrm{N}(x)\) ，这里利用了 \(\mathrm{N}(y) \in \mathrm{Ae}\) 。这就证明了 (ii)。

最后我们证明 (iii)。若 \(\mathrm{N}(x)\) 在 Ae 中可逆，且记 \(x' = \mathrm{N}(x)^{-1}\bar{x}\) ，则 \(xx' = x'x = e\) ，因为 \(\mathrm{N}(x) = x\bar{x} = \bar{x}x\) 。反之，若 x 有左逆 \(x''\) ，则 \(\mathrm{N}(x'')\mathrm{N}(x) = \mathrm{N}(e) = e\) （由 (ii)），故 \(\mathrm{N}(x)\) 在 Ae 中可逆；此外，由于 \(x' = \mathrm{N}(x)^{-1}\bar{x}\) 属于由 x 和 e 生成的子代数，元素 \(x, x', x''\) 都属于由 x、 \(x''\) 和 e 生成的结合子代数，因此 \(x'' = x''(xx') = (x''x)x' = x'\) ，唯一性由此得证。公式 \(x^{-1}(xy) = x(x^{-1}y) = y\) 由如下事实得出： \(x^{-1}\) 、x 和 y 是由 x、y 和 e 生成的子代数的元素，而该子代数是结合的。

命题 3. 设 E 是一个凯莱 A-代数， \(\gamma\) 是 A 的一个元素，F 是 E 的由 \(\gamma\) 与 E 的共轭所定义的凯莱扩张（§ 2，第 5 号，命题 5）。F 为交错代数的必要充分条件是 E 为结合代数。

设 \(u = (x, y)\) , \(v = (x', y')\) 是 \(\mathbf{F}\) 的两个元素（其中 \(x, y, x', y'\) 属于 \(\mathbf{E}\) ）。则（§ 2，第 5 号，公式 (27)）

\[
\left\{ \begin{array}{l} u ^ {2} v = ((x ^ {2} + \gamma \bar {y} y) x ^ {\prime} + \gamma \bar {y} ^ {\prime} (y \bar {x} + y x), (y \bar {x} + y x) \bar {x} ^ {\prime} + y ^ {\prime} (x ^ {2} + \gamma \bar {y} y)) \\ u (u v) = (x (x x ^ {\prime} + \gamma \bar {y} ^ {\prime} y) + \gamma (x ^ {\prime} \bar {y} + \bar {x}. \bar {y} ^ {\prime}) y, y (\bar {x} ^ {\prime} \bar {x} + \gamma \bar {y} y ^ {\prime}) + (y \bar {x} ^ {\prime} + y ^ {\prime} x) x). \end{array} \right.\tag{3}
\]

利用 \(\bar{y}y\) 和 \(\bar{x} + x\) 属于 Ae 这一事实，考察这些公式表明，E 的结合性蕴含 \(u^{2}v = u(uv)\) ，从而 F 是交错的（命题 2）。反之，若 F 是交错的，则方程 \(u^{2}v = u(uv)\) 在 \(y' = 0\) 时给出

\[
(y \bar {x} + y x) \bar {x} ^ {\prime} = y (\bar {x} ^ {\prime} \bar {x}) + (y \bar {x} ^ {\prime}) x.
\]

现在左边等于 \((y\mathrm{T}(x))\bar{x}' = y(\bar{x}'\mathrm{T}(x)) = y(\bar{x}'x + \bar{x}'\bar{x})\) ；

与右边比较，我们得到 \((y\bar{x}^{\prime})x = y(\bar{x}^{\prime}x)\) ，这证明了 E 的结合性，因为 x、y 和 \(\bar{x}^{\prime}\) 是 E 的任意元素。

